% Figure: the turbine-blade cooling trade. Raising the turbine inlet temperature
% lifts the ideal Brayton work, but the cooling air bled from the compressor to
% keep the blades below their metal limit is work removed from the cycle. The net
% specific work rises with inlet temperature at first, then bends over as the
% cooling-air fraction climbs, so a real cooled turbine has an optimum inlet
% temperature above which more cooling air costs more than the hotter gas returns.
% Two curves: the uncooled ideal (rising) and the net-of-cooling (rising then
% flattening). Illustrative shapes, not a property calculation.
\begin{figure}[htbp]
\centering
\begin{tikzpicture}
\begin{axis}[
  width=12cm, height=7.5cm,
  xlabel={turbine inlet temperature $T_3$ (K)},
  ylabel={net specific work $w_{\text{net}}$ (kJ/kg)},
  xmin=1200, xmax=2100, ymin=250, ymax=650,
  axis lines=left,
  legend pos=north west, legend cell align=left,
  tick label style={font=\scriptsize}, label style={font=\small},
  legend style={font=\scriptsize}, clip=false,
]
% Uncooled ideal: w rises roughly linearly with T3 at fixed pressure ratio.
% w_uncooled = 0.30*(T3 - 900) + 250  (illustrative, kJ/kg)
\addplot[engblue, very thick, samples=60, domain=1200:2100]
  {0.30*(x - 900) + 250};
\addlegendentry{uncooled ideal}
% Cooled net: subtract a cooling-air penalty that grows quadratically once T3
% exceeds the uncooled metal limit (~1300 K). penalty = 0 below 1300,
% then 0.0009*(T3-1300)^2 above.
\addplot[engred, very thick, samples=120, domain=1200:2100]
  {0.30*(x - 900) + 250 - (x > 1300)*0.0009*(x - 1300)^2};
\addlegendentry{net of cooling-air bleed}
% mark the optimum of the cooled curve near T3 = 1870 K (where derivative = 0)
% d/dT3[0.30 - 0.0018(T3-1300)] = 0 -> T3 = 1300 + 0.30/0.0018 = 1467? recompute
% derivative of penalty: 0.0018(T3-1300); set 0.30 = 0.0018(T3-1300) -> T3-1300=167
% -> T3 = 1467. net at 1467: 0.30*(567)+250 - 0.0009*167^2 = 170.1+250-25.1=395
\addplot[engblack, only marks, mark=*, mark size=1.8pt] coordinates {(1467,395)};
\node[engblack, font=\scriptsize, anchor=south west] at (axis cs:1467,398)
  {optimum cooled $T_3$};
% metal limit reference line
\addplot[enggray, dashed, thick, samples=2, domain=250:650] coordinates {(1300,250) (1300,650)};
\node[enggray, font=\scriptsize, anchor=south, rotate=90] at (axis cs:1300,540)
  {uncooled metal limit};
\end{axis}
\end{tikzpicture}
\caption[Turbine-blade cooling penalty]{Net specific work of a gas turbine
against turbine inlet temperature at a fixed pressure ratio. The uncooled ideal
work rises steadily with inlet temperature, but no metal survives much above the
dashed limit without cooling, so beyond it a fraction of the compressed air is
bled around the combustor to cool the first turbine stages. That bled air does no
combustion work and dilutes the gas that reaches the blades, a penalty that grows
faster than linearly as the inlet temperature and the cooling demand climb
together. The net-of-cooling curve therefore rises, bends over, and reaches an
optimum inlet temperature above which more cooling air costs more work than the
hotter gas returns. The relentless historical climb in inlet temperature has been
bought by better cooling schemes and thermal-barrier coatings that push the metal
limit and the optimum rightward, not by ignoring the penalty.}
\label{fig:vol04ch03:blade-cooling-penalty}
\end{figure}
